\documentclass[submission,copyright,creativecommons]{eptcs}
\providecommand{\event}{WPTE 2026}

\usepackage{iftex}
\ifpdf
  \usepackage{underscore}
  \usepackage[T1]{fontenc}
\else
  \usepackage{breakurl}
\fi

\usepackage{amsmath}
\usepackage{array}
\usepackage{booktabs}
\usepackage{xspace}
\raggedbottom

\title{Self-Contained Decreasing-Diagram Certificates with Canonical Compression}
\author{Anonymous
\institute{Anonymous Institution}
\email{email@example.org}
}
\def\titlerunning{Self-Contained Decreasing-Diagram Certificates}
\def\authorrunning{Anonymous}

\newcommand{\lean}{Lean\xspace}
\newcommand{\ddcert}{\texttt{ddcert}\xspace}
\newcommand{\steps}{\textit{steps}}
\newcommand{\certs}{\textit{certs}}
\newcommand{\family}[1]{\texttt{family-#1}}

\begin{document}
\maketitle

\begin{abstract}
We present a \lean-checked certificate format for finite decreasing-diagram
confluence proofs in which the exported artifact describes the entire
labeled rewrite system. A raw artifact enumerates both the labeled steps and
the peak certificates; the checker interprets the step list as the complete
underlying relation and derives confluence directly for that embedded system.
On top of the raw format we add a deterministic canonical compression layer
that removes trivial peaks and symmetric duplicates while enforcing the
resulting representation invariant during checking. The implementation comes
with positive and negative regression suites plus a generated external family
that scales through $n = 100$. For that family, compression reduces $707$
raw certificates to $101$ compressed certificates, shrinks the JSON artifact
from $115757$ bytes to $34500$ bytes, and lowers average check time from
$29.5$ ms to $18.8$ ms. The complete artifact story is reproducible from the
repository by scripts that rebuild the checker, regenerate the corpus, validate
the raw-to-canonical roundtrip, and confirm that malformed or non-canonical
artifacts are rejected.
\end{abstract}

\section{Introduction}

Decreasing diagrams are a powerful confluence criterion for abstract rewrite
systems~\cite{vanOostrom1994}. They have also been formalized in theorem
provers, including Isabelle~\cite{zankl2013}, which makes them a natural
target for certificate-oriented checking. In rewriting more broadly, the value
of exchanging proof artifacts with small trusted certifiers has already been
demonstrated by systems such as CeTA and the Certification Problem
Format~\cite{thiemann2009,sternagel2014}. Our goal is to bring that style of
artifact-oriented trust story to decreasing-diagram confluence certificates in
\lean~\cite{demoura2015}.

The key design choice is that our external certificate is \emph{self-contained}.
Instead of checking a certificate against a separately trusted rewrite relation,
the artifact itself contains the entire finite labeled relation as a list of
steps. The checker uses that list as the complete semantics of the system under
consideration. A successful check therefore certifies confluence of exactly the
relation described by the artifact, not of some external object that needs to
be reconstructed out of band.

This choice gives a clean portability story, but it also makes representation
quality matter. A raw decreasing-diagram artifact contains many certificates
that are technically valid but structurally uninteresting: trivial self-peaks
and symmetric duplicates of the same local peak. We therefore add a canonical
compressed format that retains only one orientation for each non-trivial peak,
chosen deterministically from the order of the enumerated step list. The
compressed checker is intentionally stronger than mere semantic validity: it
also enforces the canonical representation invariant. As a result, artifact size
reduction becomes a genuine technical contribution instead of an optional
post-processing step.

This paper makes four concrete contributions.
\begin{itemize}
\item We define a self-contained external certificate format for finite
  decreasing-diagram proofs and a \lean theorem that turns a successful check
  into confluence of the artifact's embedded labeled relation.
\item We define a canonical compression discipline for external artifacts and a
  strengthened checker that rejects non-canonical compressed encodings.
\item We provide a negative regression suite aimed at certificate failures that
  matter for trust: malformed JSON, missing peak coverage, bad valleys, extra
  trivial certificates, and symmetric duplicates.
\item We provide a reproducible evaluation path with both hand-written
  benchmarks and a generated family whose scaling laws are simple enough to
  analyze and large enough to be meaningful.
\end{itemize}

\section{Self-Contained Certificate Checking}

\subsection{Artifact Schema}

A raw external certificate artifact is a pair
\[
  A = (\steps(A), \certs(A))
\]
where \steps is a finite list of labeled rewrite steps and \certs is a finite
list of decreasing-diagram peak certificates. The crucial move is to derive the
rewrite relation directly from the artifact:
\[
  R_A(\ell,a,b) \iff (\ell,a,b) \in \steps(A).
\]
The checker does not appeal to any hidden trusted relation beyond this embedded
interpretation.

At the JSON level, a step is a triple of a label, a source node, and a target
node. A peak certificate records a source, the two diverging labeled steps, and
a valley certificate giving a common join together with left and right paths to
that join. The external format deliberately uses only finite first-order data:
natural numbers in the current benchmark corpus, lists, and records. This keeps
the parser small and makes the exported artifact easy to inspect independently
of the theorem prover.

The self-contained viewpoint is easiest to understand on a tiny gadget with one
non-trivial peak:
\[
  0 \xrightarrow{2} 1,\qquad
  0 \xrightarrow{2} 2,\qquad
  1 \xrightarrow{1} 3,\qquad
  2 \xrightarrow{1} 3,\qquad
  3 \xrightarrow{0} 0.
\]
In raw form, the artifact stores the five steps above, five trivial self-peak
certificates, and two non-trivial certificates for the peak at source $0$, one
for each orientation of the pair of label-$2$ steps. In compressed form, the
step list is unchanged, the trivial self-peaks are removed, and only one of the
two non-trivial peak orientations is kept.

\subsection{Checker Obligations}

Operationally, the raw self-contained checker instantiates the existing
decreasing-diagram certificate checker with $R_A$ and verifies two properties:
\begin{enumerate}
\item every enumerated step is indeed a step of $R_A$, which is immediate by
  construction; and
\item every local peak over the step list is matched by a valid certificate.
\end{enumerate}
The second condition is the substantive one. Each peak certificate names a
source, two diverging labeled steps, and a valley witnessing joinability under
labels that are smaller than both peak labels. A successful check therefore
establishes local decreasingness for all local peaks induced by the artifact.

Writing \textsf{Valid}(A,R_A,<) for the conjunction of step soundness and peak
coverage, the raw checker establishes
\[
  \texttt{selfContainedCheck}(A) = \texttt{true}
  \Longrightarrow
  \textsf{Valid}(A,R_A,<).
\]
From there the generic decreasing-diagram development supplies confluence of the
unlabeled union of $R_A$ under a well-founded label ordering. In the current
artifact corpus the label ordering is the standard order on natural numbers,
which is well founded by construction.

\paragraph{Soundness statement.}
The \lean development exports a theorem of the following form:
\[
  \texttt{selfContainedCheck}(A) = \texttt{true}
  \Longrightarrow
  \textsf{Confluent}(\bigcup R_A).
\]
The proof is short at the external layer because all the substantive
decreasing-diagram reasoning has already been established for generic labeled
ARSs. The external contribution is to instantiate those generic results with an
embedded relation whose completeness is immediate from the artifact itself.

The \lean development packages this argument into a direct confluence theorem
for self-contained artifacts: if a raw artifact passes the self-contained
checker, then the unlabeled union of $R_A$ is confluent. This is the main soundness
story for the external interface. The result is not ``the underlying system is
confluent'' for some external notion of underlying system; it is \emph{the
artifact itself denotes a confluent finite labeled ARS}.

This matters for interoperability. A consumer that receives a JSON artifact does
not need to recover the original benchmark generator, the internal
representation of the rewrite system, or a trusted parser for some other domain
language. It only needs the \ddcert checker and the artifact.

\subsection{File-Backed Regression Theorems}

The repository does not merely generate artifacts on demand. It also contains
file-backed theorems showing that selected repository artifacts are exactly the
JSON strings expected by the generator and that those file contents re-parse to
certificates accepted by the checker. This matters for the workshop setting:
the checked theorems and the exported files live in the same repository and can
be validated together.

For the hand-written external benchmarks, the repository includes checked file
regressions for the positive artifacts and for the parser/checker rejection of
negative ones. For the generated family, we currently pin file-backed checks at
\family{0} and \family{5}. This gives a small but representative bridge between
the abstract generation functions and concrete repository artifacts.

\subsection{Lean Interface}

The external layer is small enough that its main proof interface can be named
explicitly. In the raw namespace, \texttt{embeddedRelation} interprets the step
list as the whole labeled ARS, \texttt{selfContainedCheck} runs the checker
against that relation, and
\texttt{confluent\_of\_selfContainedCheck\_eq\_true} packages the final
soundness step. In the compressed namespace, \texttt{strictCheck} enforces both
semantic validity and canonical-format compliance,
\texttt{strictValid\_of\_strictCheck\_eq\_true} exposes the corresponding
logical invariant, and the compressed
\texttt{confluent\_of\_selfContainedCheck\_eq\_true} theorem again lifts a
successful self-contained check to confluence of the embedded relation.

This separation is useful operationally. A producer can emit a raw artifact
without first proving canonicity, then derive the canonical form with
\texttt{ofRaw}. A consumer that only accepts compressed artifacts can instead
rely on the stronger interface: successful checking means both confluence and
adherence to the repository's canonical representation discipline. The theorem
\texttt{strictCheck\_eq\_true\_of\_ofRaw\_valid} is the bridge between the two
views. It states that compression of a semantically valid raw artifact always
passes the stronger compressed checker.

Table~\ref{tab:interfaces} summarizes the operational distinction between the
two external interfaces.

\begin{center}
\refstepcounter{table}\label{tab:interfaces}
\small
\begin{tabular}{>{\raggedright\arraybackslash}p{1.75cm}>{\raggedright\arraybackslash}p{9.0cm}}
\toprule
Interface & Checker obligation and rejection behavior \\
\midrule
Raw & Stores all ordered local peaks, including trivial self-peaks and symmetric duplicates.
Every ordered local peak over the step list must have a valid stored certificate.
Missing self-peaks, missing symmetric orientations, and bad valleys are rejected. \\
\midrule
Compressed & Stores only non-trivial certificates in canonical orientation. Trivial
self-peaks are handled implicitly, while opposite orientations are covered through
the canonical/symmetric split. Extra trivial certificates, symmetric duplicates,
non-canonical orientations, and bad valleys are rejected. \\
\bottomrule
\end{tabular}

\smallskip
\textbf{Table~\thetable.} Raw and compressed external interfaces expose different proof obligations.
\end{center}

\section{Canonical Compression}

Raw decreasing-diagram artifacts contain two systematic sources of redundancy.
First, every step induces a trivial peak with itself; such peaks are easy to
certify but carry no structural information. Second, every non-trivial local
peak $(s_1,s_2)$ has a symmetric orientation $(s_2,s_1)$; keeping both doubles
certificate storage without adding semantic content.

Our compressed representation removes both redundancies. It stores the same
step list as the raw artifact, but it retains only those certificates that
satisfy a predicate \texttt{shouldKeep}. That predicate rejects trivial peaks
and keeps exactly one orientation of each non-trivial peak according to the
first-occurrence indices of the two steps in the artifact's step list. The step
enumeration thus induces a canonical orientation relation on local peaks.

The compressed checker is intentionally stronger than the raw checker. It does
not merely ask whether the stored certificates cover all peaks up to symmetry.
It also checks that every stored certificate is non-trivial and canonically
oriented. This strengthened validity condition gives us two benefits.

First, canonical compression is deterministic. Re-parsing a raw artifact,
compressing it, and re-serializing it yields the same compressed JSON as the
exported canonical artifact. In our generated benchmark family, the
raw-to-compressed roundtrip succeeds for every instance from $0$ to $100$.

Second, the rejection behavior becomes sharper. A compressed artifact can fail
even when its certificates would be semantically sufficient for confluence,
because the format itself has been violated. This is exactly what we want from
a certificates paper: the checker enforces not just correctness but also the
structural discipline that makes artifact exchange predictable.

\subsection{Why Canonical Orientation Matters}

It is tempting to regard canonical orientation as mere normalization of syntax,
but in our setting it carries methodological weight. The workshop claim is not
just that smaller artifacts exist. It is that there is a \emph{checker-enforced}
compressed interface whose artifacts are stable enough to exchange, regenerate,
and compare byte-for-byte.

This gives a useful separation of concerns. Raw artifacts are convenient for
producers because they can emit every peak certificate they know about, without
worrying about redundancy. Compressed artifacts are convenient for consumers
because they expose a smaller, deterministic normal form. The compressed checker
is what connects these worlds without introducing trust in an external
canonicalizer: any consumer can recompute the compression result from the raw
artifact and verify that it matches the canonical compressed JSON.

\subsection{Strict Validity}

The strengthened compressed checker can be described as a \emph{strict validity}
test:
\[
  \textsf{StrictValid}(A,R_A,<) \equiv
  \textsf{Valid}(A,R_A,<) \wedge \textsf{CanonicalCerts}(A).
\]
The first conjunct is semantic validity up to the compressed peak-coverage
criterion; the second says that every stored certificate satisfies the
canonical-keep predicate induced by the step enumeration. The external layer
proves both directions needed for use:
\begin{itemize}
\item successful compressed checking implies strict validity; and
\item compression of a valid raw artifact produces a strictly valid compressed
  artifact.
\end{itemize}
The second point is what powers the roundtrip validation experiment described
below.

\section{Evaluation and Regressions}

\subsection{Benchmarks}

Our evaluation combines seven hand-written self-contained external benchmarks
with a generated external family. The hand-written suite now covers wider
branching, empty-path shortcuts, longer asymmetric valleys, mixed three-way
peaks, and multiple independent peaks. The family instance \family{n}
consists of $n+1$ independent peak gadgets. This yields simple exact counts:
\[
\begin{aligned}
  |\steps(\family{n})| &= 5(n+1), \\
  |\certs_{\mathrm{raw}}(\family{n})| &= 7(n+1), \\
  |\certs_{\mathrm{comp}}(\family{n})| &= n+1.
\end{aligned}
\]
The generated family is therefore not just a source of large inputs; it gives a
transparent scaling law for how much redundancy compression removes.

\subsection{Generated Family Construction}

Each generated gadget has a hub node, a left branch, a right branch, a join
node, and a loop back to the hub. The two outgoing hub steps carry the same peak
label, the two downward branch steps carry a smaller label, and the loop step
uses label $0$. This design gives exactly one non-trivial local peak per gadget,
with a short valley on each side. Because the gadgets are independent, the
global benchmark family is easy to analyze exactly, yet it still exercises the
same parser, checker, JSON renderer, and canonicalizer as the hand-written
benchmarks.

This family is a better workshop benchmark than a grab bag of unrelated random
examples. It demonstrates that compression behavior is structural, not an
accident of a few cases: every instance has the same 85\% certificate reduction,
and the JSON savings stabilize around 70\% across the entire measured range.
Likewise, the raw and compressed check times grow smoothly with the number of
gadgets, while compressed checking remains strictly faster than raw checking in
every reported instance.

Table~\ref{tab:benchmarks} reports representative artifact sizes and average
check times from the repository's regenerated CSV summaries. Compression is
consistently beneficial: it reduces both artifact size and checking time on all
measured cases. For \family{100}, the reduction is from $707$ raw certificates
to $101$ compressed certificates and from $115757$ bytes to $34500$ bytes,
while average check time drops from $29501$ $\mu$s to $18766$ $\mu$s.

\begin{table}[t]
\centering
\small
\begin{tabular}{lrrrrrrrr}
\toprule
Case & Steps & Raw c. & Comp. c. & Red. & Raw B & Comp. B & Raw $\mu$s & Comp. $\mu$s \\
\midrule
\texttt{fan3} & 7 & 13 & 3 & 76\% & 1905 & 655 & 66 & 30 \\
\texttt{staggered} & 6 & 8 & 1 & 87\% & 1233 & 363 & 41 & 15 \\
\texttt{double-peak} & 9 & 13 & 2 & 84\% & 1953 & 591 & 66 & 27 \\
\family{0} & 5 & 7 & 1 & 85\% & 1065 & 323 & 35 & 13 \\
\family{5} & 30 & 42 & 6 & 85\% & 6465 & 1902 & 273 & 123 \\
\family{20} & 105 & 147 & 21 & 85\% & 23055 & 6792 & 1763 & 970 \\
\family{50} & 255 & 357 & 51 & 85\% & 57357 & 17000 & 8258 & 5034 \\
\family{100} & 505 & 707 & 101 & 85\% & 115757 & 34500 & 29501 & 18766 \\
\bottomrule
\end{tabular}
\caption{Representative external benchmark data from the regenerated artifact corpus.}
\label{tab:benchmarks}
\end{table}

The hand-written benchmarks and generated family play complementary roles.
The former show that the format is expressive enough for bespoke examples with
different local-peak shapes, including larger fans, shortcut valleys, and
longer asymmetric joins. The latter show that the same interface scales in a
predictable way and that the canonical compressed format has a clear asymptotic
story instead of isolated wins.

\subsection{Negative Regressions}

Positive examples are not enough for a certificates paper. We also need to show
that the checker rejects the kinds of artifacts that would otherwise undermine
the trust story. Table~\ref{tab:negative} summarizes the current negative
regression corpus.

\begin{table}[t]
\centering
\small
\begin{tabular}{p{5.3cm}p{7.0cm}p{1.6cm}}
\toprule
Artifact & Failure mode & Outcome \\
\midrule
\texttt{fan3-missing-loop-self.raw} & missing raw self-peak certificate for the loop step & reject \\
\texttt{fan3-extra-trivial} & extra trivial certificate in compressed form & reject \\
\texttt{fan3-missing-peak} & incomplete peak coverage & reject \\
\texttt{staggered-bad-valley} & invalid valley witness & reject \\
\texttt{malformed} & malformed JSON input & parse fail \\
\texttt{family-5-missing-self.raw} & generated raw artifact with missing loop self-peak & reject \\
\texttt{family-5-extra-trivial} & non-canonical extra trivial certificate & reject \\
\texttt{family-5-symmetric-duplicate} & symmetric duplicate survives compression & reject \\
\texttt{family-5-missing-peak} & generated compressed artifact with one peak removed & reject \\
\texttt{family-5-bad-valley} & generated compressed artifact with invalid valley endpoint & reject \\
\bottomrule
\end{tabular}
\caption{Negative regressions exercised by the current checker.}
\label{tab:negative}
\end{table}

These regressions are not decorative. The last four lines are specifically about
the canonical compressed interface, not just semantic confluence checking.
Without them, it would be hard to argue that the compressed format contributes
anything stronger than a convenient file-size reduction.

\subsection{Roundtrip Validation}

The generated family also validates the determinism of canonical compression.
For every instance between $0$ and $100$, the command
\texttt{roundtrip-external-family-range} reparses the raw artifact, recomputes
the compressed form, and checks that the serialized result matches the exported
canonical JSON exactly. The resulting CSV reports \texttt{ok} for all $101$
instances. This is a compact way to demonstrate that canonicalization is not an
informal convention but a stable part of the checker contract.

\section{Implementation Structure}

The implementation is split across a small number of layers. The generic
decreasing-diagram certificate checker lives in the core rewriting
infrastructure. The external self-contained layer introduces the embedded
relation semantics and the stronger compressed checker. On top of that sit two
benchmark suites: a hand-written suite of positive and negative examples, and
the generated external family. Finally, the \ddcert command-line tool provides
artifact export, checking, compression, statistics, and roundtrip commands.

This stratification is useful for two reasons. First, it keeps the trusted story
simple. The external layer does not re-prove the generic decreasing-diagram
machinery; it reuses it through small instantiation theorems. Second, it keeps
evaluation code honest. The same exported commands that create the benchmark
tables are the ones exercised by the regeneration scripts, so the paper is
reporting on the interface that users would actually invoke.

\subsection{Artifact Bundles and Scripts}

The repository currently exports four artifact bundles relevant to the workshop
paper:
\begin{itemize}
\item the small built-in example and typed internal family;
\item the hand-written external positive benchmarks;
\item the hand-written and generated negative external artifacts; and
\item the generated external family with both raw and compressed JSON plus CSV
  summaries.
\end{itemize}
The repository includes a general regeneration script for the decreasing-diagram
bundles. A separate workshop-preparation script runs the full build,
regenerates the artifacts at a larger range, and copies the resulting CSV files
and summary text into a dedicated workshop directory.

\subsection{What Is Checked in Lean}

Several claims in the evaluation are reflected directly by \lean theorems rather
than only by script output. These include successful self-contained checking for
sampled positive family instances, rejection of sampled negative generated
instances, confluence of selected file-backed compressed artifacts, and equality
between selected file contents and the strings produced by the generator. The
performance tables themselves come from the CLI benchmarking code, but the
semantic correctness claims are anchored in checked theorems.

\section{Discussion}

\subsection{Positioning Relative to Existing Certificate Ecosystems}

Our work is close in spirit to certificate-oriented rewriting infrastructure
such as CPF and CeTA~\cite{sternagel2014,thiemann2009}, but it makes a
different design tradeoff. Those ecosystems aim to cover a broad space of proof
objects and external tools. By contrast, our current external interface is
narrower and more self-contained: a certificate names its entire finite labeled
relation explicitly. This makes the trust boundary particularly simple, at the
cost of artifact sizes that scale with the explicit step list.

Relative to the prior formalization of confluence by decreasing diagrams in
Isabelle~\cite{zankl2013}, the novelty here is not the criterion itself. The
novelty is the certificate interface around it: external JSON artifacts,
embedded-relation semantics, canonical compressed checking, and a reproducible
artifact corpus. In other words, we are pushing decreasing diagrams from a
formalized proof principle toward a portable certificate technology.

\subsection{Trusted Computing Base and Threat Model}

The trust story is intentionally narrow. The producer of a certificate is
untrusted: it may omit peaks, invent valleys, emit malformed JSON, or output a
compressed artifact that is semantically adequate but not canonical. The
consumer trusts only the checker stack used to parse and validate the artifact.
At proof level this reduces to the \lean kernel together with the imported
decreasing-diagram development. At executable level it also includes the
compiled JSON parser and runtime used by \ddcert.

The self-contained step list shrinks this boundary in an important way. The
consumer does not need to trust a separate benchmark generator, a parser for an
external rewriting language, or an out-of-band reconstruction of the rewrite
relation. All such producers can be treated as arbitrary artifact emitters. If
their output is inconsistent with the checker contract, the artifact is
rejected.

Our negative regression corpus is organized around this threat model. Missing
self-peaks and missing non-trivial peaks test coverage failures. Bad valleys
test semantic unsoundness. Malformed JSON tests the parser boundary. Extra
trivial certificates and symmetric duplicates test attempts to smuggle
non-canonical compressed artifacts past the consumer. File-backed theorems
address a different but practical risk: drift between generator code and
committed artifact files. Together, these checks do not verify the operating
system, compiler, or runtime environment, but they do substantially reduce the
amount of domain-specific trust required to exchange decreasing-diagram
certificates.

\subsection{Limits}

The current external interface is intentionally finite and first-order. It is
well suited to compact confluence certificates for small and medium finite
systems, but it is not yet a general-purpose interchange format for arbitrary
rewriting tools. We also currently benchmark deterministically generated
families rather than certificates emitted by independent external provers.

These are reasonable limitations for a workshop paper. The present contribution
is a crisp interface with a clean trust story and a reproducible implementation.
The obvious next step is integration: emitting these certificates from external
confluence tools or relating the format more directly to existing rewriting
certificate standards.

\section{Reproducibility}

The repository contains the entire artifact pipeline. Running
\texttt{scripts/prepare-wpte-2026.sh 20 100} performs a full
\texttt{lake build}, regenerates the hand-written and generated decreasing-diagram
artifacts, refreshes the benchmark CSV files, checks the largest positive
generated artifacts in both raw and compressed form, validates the
raw-to-canonical roundtrip through \family{100}, confirms that the generated
negative family artifacts are rejected, and copies the main evaluation outputs
into the directory \texttt{artifacts/decreasing-diagrams/wpte-2026/}.

This matters because the paper's technical claims line up directly with
repository outputs:
\begin{itemize}
\item self-contained checking is exercised by the external raw and compressed
  checker commands of \ddcert;
\item canonical compression is exercised by the compression command and by the
  generated-family roundtrip validation;
\item positive evaluation is summarized by the hand-written benchmark CSV and
  the generated-family statistics CSV; and
\item negative evaluation is witnessed by the rejection of the generated
  family-negative artifact suite.
\end{itemize}
As a result, the submission does not depend on a hidden benchmarking script or a
manually assembled appendix. The paper and artifact story are generated from the
same checked implementation.

\section{Conclusion}

We described a self-contained external certificate interface for
decreasing-diagram confluence proofs, implemented in \lean and backed by a
canonical compressed representation. The main point is not merely that the
artifacts can be made smaller. Rather, the combination of an embedded relation,
strongly checked canonical compression, negative regressions, and reproducible
benchmark generation yields a cleaner trust story for certificate exchange.

There are at least two natural next steps. The first is to connect this
self-contained interface to external confluence tools, so that automatically
generated certificates can be exported directly in canonical form. The second is
to compare the design more explicitly with existing rewriting certificate
ecosystems, especially CPF-style exchange formats, and to investigate whether
self-contained decreasing-diagram certificates can serve as one component in a
broader certified confluence pipeline.

\begin{thebibliography}{10}

\bibitem{demoura2015}
Leonardo de Moura, Soonho Kong, Jeremy Avigad, Floris van Doorn \& Jakob von
Raumer (2015):
\newblock The Lean Theorem Prover (System Description).
\newblock In Amy P. Felty \& Aart Middeldorp, editors:
\newblock {\em Automated Deduction -- CADE-25},
\newblock LNCS 9195, Springer, pp.\ 378--388.

\bibitem{sternagel2014}
Christian Sternagel \& Ren{\'e} Thiemann (2014):
\newblock The Certification Problem Format.
\newblock {\em arXiv}:1410.8220.

\bibitem{thiemann2009}
Ren{\'e} Thiemann \& Christian Sternagel (2009):
\newblock Certification of Termination Proofs Using CeTA.
\newblock In Stefan Berghofer, Tobias Nipkow, Christian Urban \& Makarius
Wenzel, editors:
\newblock {\em Theorem Proving in Higher Order Logics},
\newblock LNCS 5674, Springer, pp.\ 452--468.

\bibitem{vanOostrom1994}
Rob van Oostrom (1994):
\newblock Confluence by Decreasing Diagrams.
\newblock {\em Theoretical Computer Science} 126(2), pp.\ 259--280.

\bibitem{zankl2013}
Harald Zankl (2013):
\newblock Confluence by Decreasing Diagrams -- Formalized.
\newblock In Femke van Raamsdonk, editor:
\newblock {\em 24th International Conference on Rewriting Techniques and
Applications (RTA 2013)},
\newblock LIPIcs 21, Schloss Dagstuhl -- Leibniz-Zentrum f{\"u}r Informatik,
pp.\ 352--367.

\end{thebibliography}

\end{document}
